\section{Introduction}

Long-horizon agentic tasks require many turns of Large Language Model (LLM) interaction with an environment. This includes tasks such as conversational tool use \citep{schick2023toolformer, qin2023toolllm}, agentic coding \citep{jimenez2024swebench}, terminal interaction \citep{xie2024osworld}, and web search \citep{yao2022webshop, wei2025browsecompsimplechallengingbenchmark}. These tasks have emerged as the frontier in AI systems capable of executing complex, real-world workflows. However, post-training for these long-horizon agentic capabilities introduces a fundamental tension between training efficiency and generalization. 
While supervised fine-tuning (SFT) offers a compute-efficient mechanism for acquiring new capabilities, it frequently struggles to generalize beyond the training distribution and catastrophically degrades out-of-domain (OOD) performance \citep{chu2025sftmemorizesrlgeneralizes, luo2023empiricalCF}. In contrast, end-to-end reinforcement learning (E2E RL) typically yields higher in-domain accuracy while robustly retaining OOD capabilities \citep{ouyang2022training, chen2025retainingdoingroleonpolicy}. Nevertheless, E2E RL incurs high compute overheads, as each parameter update necessitates repeated, many turns of on-policy rollout with environmental interactions. This dichotomy naturally raises the following question:
\begin{center}
    \emph{Can we combine the data efficiency of SFT with the generalization capabilities of E2E RL, achieving both in-domain accuracy and OOD retention without incurring full-trajectory rollouts?}
\end{center}
To address this challenge, a natural {attempt would be to repurpose SFT trajectory demonstrations for RL.} 
% baseline would repurpose SFT trajectories for RL. 
Specifically, one could sample on-policy rollouts conditioned on a randomly selected intermediate turn from an SFT trajectory, assigning a positive reward only if the generated action exactly matches the SFT data demonstration. While this mitigates the under-utilization of SFT data which only trains on a single demonstrated completion, our preliminary experiments reveal that it fails to improve OOD accuracy relative to standard SFT on the same data. We empirically trace this failure to two bottlenecks. First, randomly selected intermediate turns frequently provide negligible learning signals; under Group Relative Policy Optimization (GRPO)~\citep{shao2024deepseekmath}, sampled actions at such turns often uniformly succeed or fail, yielding a normalized advantage close to zero and thus producing no meaningful gradient update. Second, exact string matching with SFT data is excessively restrictive in generative action spaces, where numerous functionally equivalent actions validly diverge from the singular demonstrated completion. For instance, a tool call or search query may be perfectly appropriate despite not perfectly matching the demonstration. 

Motivated by these insights, we introduce \textbf{PivotRL}, a novel framework for long-horizon agentic RL to overcome both bottlenecks simultaneously. PivotRL executes local, on-policy rollouts and filters for \emph{pivots}: informative assistant turns where sampled actions exhibit mixed outcomes (i.e., both successes and failures). Consequently, training requires only brief, partial rollouts from these pivot states rather than exhaustive full trajectories. To effectively score these rollouts, we leverage domain-appropriate verifiers to assign rewards for functionally equivalent actions rather than strictly penalizing deviations from SFT data demonstrations \citep{shao2024deepseekmath, rohatgi2025tamingprocessverifiers}.

We substantiate our methodology with a lightweight theoretical analysis on our core design choices. 
{First, we prove that the Fisher norm of the natural gradient of the statewise reward objective scales with the reward standard deviation. Consequently, the GRPO update along the KL path scales directly with this variance, validating our strategy of filtering for mixed-outcome pivots to maximize the local in-domain learning signals.}
{Second, we show that functional reward-based RL shifts probability mass toward actions that are functionally equivalent to the expert demonstration, while preserving the conditional distributions on all other actions. This maintains the reference policy's relative ordering of task-unrelated actions, thereby mitigating out-of-domain (OOD) degradation.} 
% First, a KL-projection theorem demonstrates that the functional reward induces a conservatively regularized update relative to the reference policy: it optimally rescales the probability mass of acceptable versus unacceptable actions while strictly preserving the reference policy's relative action ordering within each respective set. 
% Second, we prove that the Fisher norm of the natural gradient of the statewise reward objective exactly equals the reward standard deviation. Consequently, the GRPO update along the KL path scales directly with this variance, mathematically validating our strategy of filtering for mixed-outcome pivots to maximize the local learning signal.

Empirically, we first show that under identical training data---the same prompts and expert trajectories---PivotRL yields larger in-domain improvements than SFT while exhibiting substantially less OOD regression. 
% Across four diverse agentic benchmarks, 
{When training in the conversational tool use, agentic coding, terminal tool use and search domain,}
PivotRL achieves an average in-domain accuracy improvement of $+14.11$\footnote{Throughout this paper, performance scores and differences are reported in percentage points.}over the base model, which significantly improves the $+9.94$ points attained by SFT. Crucially, PivotRL does not introduce OOD regressions (+$0.21$ change), whereas SFT suffers a severe regression of $-9.48$ in non-agentic domains including math, science QA and competitive coding. Second, we provide a direct comparison with E2E RL on SWE-Bench, a rigorous standard for evaluating long-horizon agentic capabilities with well-established E2E RL baselines~\citep{jimenez2024swebench, wang2025openhandsopenplatformai}. On this benchmark, we demonstrate that PivotRL attains competitive accuracy with E2E RL, but only with a {$4\times$ fewer rollout turns.} 
{PivotRL is deployed at production scale for leading open LLMs. Along with SFT and E2E RL, it serves as the primary workhorse for agentic post-training for NVIDIA's Nemotron-3-Super LLM~\citep{nvidia2025nemotron3super}.}
% Finally, rigorous ablations of our framework confirm that both informative pivot selection and functional reward design are indispensable for realizing optimal performance.

% Beyond its empirical performance, PivotRL demonstrates  scalability in real world production. Along with SFT and E2E RL, PivotRL served as the primary workhorse for aligning the majority of agentic domains during the post-training of NVIDIA-Nemotron-3-Super-120B-A12B-BF16~\citep{nvidia2025nemotron3super}. 